% ============================================================
%  EXAMEN TEMA 3: POTENCIACIÓN DE NÚMEROS ENTEROS
%  Curso Introductorio de Matemáticas — 1er Año
%  Autor: Ing. Gabriel Astudillo
%  Sesiones cubiertas: Semana 3 (clases 9 a 12)
%  Nota: enunciado limpio (sin pistas ni desarrollo). Las
%  soluciones están en la "Clave de Respuestas" al final.
% ============================================================

\documentclass[12pt, letterpaper]{article}

% ── Paquetes ─────────────────────────────────────────────────

\usepackage{mathwizards-guia}
\mwguia{Examen — Tema 3: Potenciación}{Ing. Gabriel Astudillo $\cdot$ Curso 1er Año}

% ── Comandos propios de este examen ───────────────────
\newcommand{\Z}{\mathbb{Z}}

\begin{document}
% ============================================================

% ── Portada / Título ─────────────────────────────────────────
\mwportada{Examen del Tema 3}{Potenciación de Números Enteros}{Definición, Propiedades y Primeros Pasos en Álgebra}

\vspace{4pt}
\begin{cuadroinfo}
  \small
  \textbf{Nombre:} \rule{0.55\linewidth}{0.4pt} \quad \textbf{Fecha:} \rule{0.15\linewidth}{0.4pt} \\[6pt]
  \textbf{Tiempo sugerido:} 60 minutos \quad$\bullet$\quad \textbf{Puntaje total:} 100 puntos \\[4pt]
  \textbf{Instrucciones:}
  \begin{enumerate}[leftmargin=*]
    \item Lee cada ejercicio con calma antes de resolverlo.
    \item Muestra \emph{todos} tus pasos: se evalúa cómo llegas a la respuesta, no solo la respuesta.
    \item \textbf{Ojo con el signo}: recuerda la diferencia entre $(-a)^n$ y $-a^n$.
    \item Aplica correctamente las \emph{propiedades de las potencias} (producto, cociente, potencia de potencia).
    \item No se permite el uso de calculadora.
  \end{enumerate}
\end{cuadroinfo}

\vspace{8pt}

% ============================================================
\section{Base, Exponente y Signo del Resultado}
% ============================================================

\begin{enumerate}[leftmargin=*]
  \item Calcule, distinguiendo el papel del paréntesis:
        \[
          (-2)^4, \quad -2^4, \quad (-3)^3, \quad -3^3
        \]

  \item Calcule el valor exacto de:
        \[
          (-1)^{2025} + (-1)^{2026} + 0^{7} + 10^{0}
        \]

  \item Determine si cada resultado es positivo, negativo o cero, y justifique según la paridad del exponente:
        \[
          (-7)^{12}, \quad (-8)^{9}, \quad (-5)^{0}, \quad (-2)^{31}
        \]

  \item Complete las potencias de base $2$ y $3$ y úselas para calcular:
        \[
          2^{8} - 3^{3} + 2^{5} - 3^{2}
        \]
\end{enumerate}

\newpage

% ============================================================
\section{Propiedades de las Potencias}
% ============================================================

\begin{enumerate}[leftmargin=*]
  \item Simplifique aplicando la propiedad del producto de igual base:
        \[
          2^{6} \times 2^{4} \times 2^{7}
        \]

  \item Simplifique aplicando la propiedad del cociente de igual base:
        \[
          \frac{9^{8} \times 9^{3}}{9^{6}}
        \]

  \item Simplifique aplicando la propiedad de la potencia de una potencia:
        \[
          \big[(2^{3})^{4}\big]^{2}
        \]

  \item Calcule combinando varias propiedades:
        \[
          5^{4} \times 5^{2} \div 5^{3} \times 5^{0}
        \]
\end{enumerate}

\newpage

% ============================================================
\section{Potencias con Letras (Sustitución y Simplificación)}
% ============================================================

\begin{enumerate}[leftmargin=*]
  \item Sustituya $a = 3$ y $b = 2$ en la expresión:
        \[
          2a^{2}b^{3} - 3a^{2}b^{2}
        \]

  \item Con $x = 5$, calcule:
        \[
          x^{3} + 2x^{2} - 4x + 7
        \]

  \item Simplifique aplicando las propiedades:
        \[
          (m^{3a} \cdot m^{7a})^{3}
        \]

  \item Simplifique:
        \[
          \big[(3x)^{2} \cdot (2x^{3})^{2}\big]^{3}
        \]
\end{enumerate}

\newpage

% ============================================================
\section{Exponentes Negativos, Cocientes y Cambio de Base}
% ============================================================

\begin{enumerate}[leftmargin=*]
  \item Simplifique la expresión:
        \[
          \left(\frac{a^{6x}}{a^{2x}}\right)^{-2}
        \]

  \item Calcule, recordando que un exponente negativo invierte la base:
        \[
          2^{-3} + 3^{-2} - \left(\frac{1}{2}\right)^{-1}
        \]

  \item Simplifique escribiendo todo en base $2$:
        \[
          (64^{5x} \div 128^{x})^{5}
        \]

  \item Simplifique escribiendo todo en base $3$:
        \[
          (27^{p} \cdot 9^{3p})^{6}
        \]
\end{enumerate}

\newpage

% ============================================================
\ifmwclaves
  \section{Clave de Respuestas}
  % ============================================================

  \subsection*{Base, Exponente y Signo del Resultado (Ejercicios 1--4)}

  \begin{enumerate}[leftmargin=*, label=\textbf{\arabic*.}]
    \item $(-2)^4 = 16$; \ $-2^4 = -16$; \ $(-3)^3 = -27$; \ $-3^3 = -27$
    \item $(-1)^{2025} = -1$; $(-1)^{2026} = 1$; $0^{7} = 0$; $10^{0} = 1$ \ $\Rightarrow$ $-1 + 1 + 0 + 1 = 1$
    \item $(-7)^{12} > 0$ (exponente par); \ $(-8)^{9} < 0$ (exponente impar);
          $(-5)^{0} = 1 > 0$; \ $(-2)^{31} < 0$ (exponente impar)
    \item $2^{8} = 256$; $3^{3} = 27$; $2^{5} = 32$; $3^{2} = 9$ \ $\Rightarrow$ $256 - 27 + 32 - 9 = 252$
  \end{enumerate}

  \subsection*{Propiedades de las Potencias (Ejercicios 1--4)}

  \begin{enumerate}[leftmargin=*, label=\textbf{\arabic*.}]
    \item $2^{6+4+7} = 2^{17}$
    \item $\dfrac{9^{8+3}}{9^{6}} = 9^{11-6} = 9^{5}$
    \item $2^{3 \times 4 \times 2} = 2^{24}$
    \item $5^{4+2-3+0} = 5^{3} = 125$
  \end{enumerate}

  \subsection*{Potencias con Letras (Ejercicios 1--4)}

  \begin{enumerate}[leftmargin=*, label=\textbf{\arabic*.}]
    \item $2 \times 3^{2} \times 2^{3} - 3 \times 3^{2} \times 2^{2}$ \ $\Rightarrow$ $2 \times 9 \times 8 - 3 \times 9 \times 4 = 144 - 108 = 36$
    \item $5^{3} + 2 \times 5^{2} - 4 \times 5 + 7 = 125 + 50 - 20 + 7 = 162$
    \item $(m^{3a + 7a})^{3} = (m^{10a})^{3} = m^{30a}$
    \item $(3x)^{2} = 9x^{2}$; $(2x^{3})^{2} = 4x^{6}$; producto $= 36x^{8}$; $\Rightarrow (36x^{8})^{3} = 36^{3}x^{24} = 46656\,x^{24}$
  \end{enumerate}

  \subsection*{Exponentes Negativos, Cocientes y Cambio de Base (Ejercicios 1--4)}

  \begin{enumerate}[leftmargin=*, label=\textbf{\arabic*.}]
    \item $\left(\dfrac{a^{6x}}{a^{2x}}\right)^{-2} = (a^{4x})^{-2} = a^{-8x} = \dfrac{1}{a^{8x}}$
    \item $2^{-3} = \frac{1}{8}$; $3^{-2} = \frac{1}{9}$; $\left(\frac{1}{2}\right)^{-1} = 2$; \ $\frac{1}{8} + \frac{1}{9} - 2 = \frac{17}{72} - 2 = -\frac{127}{72}$
    \item $64 = 2^{6}$; $128 = 2^{7}$ \ $\Rightarrow$ $(2^{6 \cdot 5x} \div 2^{7x})^{5} = (2^{30x} \div 2^{7x})^{5} = (2^{23x})^{5} = 2^{115x}$
    \item $27 = 3^{3}$; $9 = 3^{2}$ \ $\Rightarrow$ $(3^{3p} \cdot 3^{6p})^{6} = (3^{9p})^{6} = 3^{54p}$
  \end{enumerate}

  \vspace{10pt}
  \begin{cuadrosolucion}
    \textbf{\textquestiondown Cómo te fue?} El secreto de este tema está en el \emph{signo} (paréntesis vs. sin paréntesis) y en aplicar las \emph{propiedades} con calma. Si lo dominas, \textquestiondown listo para el Tema 4!
  \end{cuadrosolucion}

\fi
\end{document}
